\section{}
\noindent \textbf{Assignment: 01 - b}
\hfill Date: 19/01/2026 \\
\noindent
\textbf{Title:} Implementation of secure socket communication using RSA asymmetric cipher.

\subsection*{Code}
\begin{lstlisting}[style=mypython, caption={TerminalMessage.py}, label={code:A-1-b-TerminalMessage}]
import colorama as CM

terminalIcon = "[>]"

def PrintTry(status:bool, message:str):
    '''
        TODO: Function description
    '''

    if(status):
        statusMessage = "Success"
        color = CM.Fore.GREEN
    else:
        statusMessage = "Error"
        color = CM.Fore.RED

    print(color + f"{terminalIcon} {statusMessage}: {message}")

def PrintMessage(message:str):
    '''
        TODO: Function description
    '''

    print(CM.Fore.MAGENTA + f"{terminalIcon} {message}")
\end{lstlisting}

\begin{lstlisting}[style=mypython, caption={Encryption.py}, label={code:A-1-b-Encryption}]
from Crypto.PublicKey import RSA
from Crypto.Cipher import PKCS1_OAEP

class rsa:
    def generateKeys(self):
        key = RSA.generate(1024)
        privateKey = key
        publicKey = key.publickey()
        return(publicKey, privateKey)

    def encrypt(self, data, publicKey):
        cipher = PKCS1_OAEP.new(publicKey)
        ciphertext = cipher.encrypt(data.encode())
        return ciphertext

    def decrypt(self, data, privateKey):
        cipher = PKCS1_OAEP.new(privateKey)
        plaintext = cipher.decrypt(data)
        return plaintext.decode("utf-8")
\end{lstlisting}

\begin{lstlisting}[style=mypython, caption={server.py}, label={code:A-1-b-server}]
import socket
import TerminalMessage as TM
from Crypto.PublicKey import RSA
from Crypto.Cipher import PKCS1_OAEP

# Server configuration.
serverConfig = {
    "name": "NITJSR Network Security Lab",
    "ip": "localhost",
    "port": 8080
}

# TODO: Create socket.
try:
    server = socket.socket()
    TM.PrintTry(True, "Scoket created successfully.")
except:
    TM.PrintTry(False, "Error: Unable to create socket.")
    exit()

# TODO: Bind the socket using the serverConfig.
bindingConfig = (serverConfig["ip"], serverConfig["port"])
try:
    server.bind(bindingConfig)
    TM.PrintTry(True, f"Success: Server binded to {bindingConfig}")
except:
    TM.PrintTry(False, "Error: Unable to bind socket.")
    exit()

# TODO: Config the max no of connections it can accept and start listening.
noOfConnections = 5
try:
    server.listen(noOfConnections)
    TM.PrintTry(True, f"Success: Server is in listening mode with max connections: {noOfConnections}")
except:
    TM.PrintTry(False, "Error: Unable to start listening.")
    exit()

'''
    TODO: Run a loop that will stop when we interupt it or an error occured, inside the loop:
            1. Accept the new connection
            2. Send connection successful message
            3. Generate RSA keys
            4. Send public key to client
            5. Ask for his name and respond with "Welcome {name}"
'''
while True:
    TM.PrintMessage("Waiting for an incomming connection.")

    # Accept the incomming connection.
    client, address = server.accept()
    TM.PrintMessage(f"Incomming connection from {address}")
    TM.PrintMessage(f"Connection accepted from {address}")

    # Send connection successful response.    
    connectionSuccessfulMessage = "Success: Connected to server.".encode()
    client.send(connectionSuccessfulMessage)

    # TODO: Generate RSA keys.
    key = RSA.generate(1024)
    privateKey = key
    publicKey = key.publickey()
    rawPublicKey = publicKey.exportKey()
    cipher = PKCS1_OAEP.new(privateKey)

    # TODO: Send public key to the client.
    client.send(rawPublicKey)   

    # Ask for name.
    message = "Please provide you name.".encode()
    client.send(message)
    
    # Accept the name.
    name = client.recv(1024)

    # TODO: Decrypt the name.
    decryptedName = cipher.decrypt(name).decode()

    # Send welcome message.
    welcomeMessage = f'{decryptedName}, Welcome to {serverConfig["name"]}\'s server.'.encode()
    client.send(welcomeMessage)
    
    # Terminate the connection.
    client.close()
    TM.PrintMessage(f"Connection terminated from {address}")
\end{lstlisting}

\begin{lstlisting}[style=mypython, caption={client.py}, label={code:A-1-b-client}]
import socket
import TerminalMessage as TM
from Crypto.PublicKey import RSA
from Crypto.Cipher import PKCS1_OAEP

# TODO: Create socket.
client = socket.socket()

# Server configuration.
config = {
    "ip": "localhost",
    "port": 8080
}
serverConfig = (config["ip"], config["port"])

# TODO: Connect to server.
try:
    client.connect(serverConfig)
    
    # TODO: Receive the connection successful message.
    connectionSuccessfulMessage = client.recv(1024).decode()
    TM.PrintTry(True, connectionSuccessfulMessage)

    # TODO: Receive the RSA public key.
    receivedPublicKey = client.recv(1024)
    publicKey = RSA.import_key(receivedPublicKey)
    TM.PrintTry(True, "Received public key for RSA")
except:
    TM.PrintTry(False, "Unable to connect to server.")
    exit()

# Create object for RSA.
cipher = PKCS1_OAEP.new(publicKey)

# TODO: Receive the ask name message and print.
nameMessage = client.recv(1024).decode()
TM.PrintMessage(nameMessage)

# TODO: Take name input and encrypt the name.
TM.PrintMessage("Your Name Please: ")
name = input()
encryptedName = cipher.encrypt(name.encode())

# TODO: Send encrypted name to the server.
client.send(encryptedName)

# TODO: Print the welcome message.
welcomeMessage = client.recv(1024)
TM.PrintMessage(welcomeMessage.decode())

# TODO: Terminate the connection.
client.close()
\end{lstlisting}

\subsection*{Output}
\begin{figure}[h]
    \centering
    \includegraphics[width=1\textwidth]{assets/img/experiment-02/output.png}
    \caption{Secure socket communication b/w server and client using RSA asymmetric cipher.}
    \label{fig:A-1-b-output}
\end{figure}

\newpage